\chapter{Conclusion and Future Work}\label{ch:conclusion}

This chapter summarizes the main results of this thesis and discusses possible future work. The conclusions are organized around the receiver-side factors studied in the previous chapters, including channel information, decoder reliability, frequency-domain diversity, and synchronization conditions in uplink NOMA systems.

\section{Conclusion}

This thesis draft compared uplink NOMA receivers from uncoded single-resource systems to coded multi-resource and OFDM systems. In the single-resource case, GDMA outperforms PD-NOMA with soft SIC because GDMA jointly exploits complex channel gains, while SIC depends heavily on received-power ordering and suffers from error propagation. Increasing the large-scale fading disparity helps SIC, but it also creates user-fairness concerns.

In the uncoded multi-resource case, MMSE-SIC improves PD-NOMA by using vector observations across resources. Post-SINR ordering is more effective than amplitude-based ordering because it reflects post-filtering interference and noise enhancement. Nevertheless, HDS-CDMA with ML detection still achieves better BER because it performs joint multiuser detection.

In coded systems, MMSE-SIC-IDD becomes much more competitive. Decoder-aided soft cancellation reduces SIC error propagation, and variance-aware MMSE filtering uses decoder reliability to avoid overconfident cancellation. MMSE-PIC-IDD and MAP-IDD provide useful comparisons: PIC reduces sequential latency, while MAP offers a high-performance but high-complexity benchmark. The coded results show that receiver architecture and channel coding should be considered jointly rather than separately, because decoder feedback changes the quality of the cancellation signal.

For coded OFDM NOMA, channel selectivity, detection ordering, and the choice between SIC and PIC all affect BLER. The OFDM extension shows that the same MMSE-SIC-IDD principle can be applied to frequency-domain observations, but coded-block ordering and subcarrier reliability become additional design issues. In asynchronous OFDM NOMA, timing offsets within the CP introduce effective channel phase rotations and possible additional interference; under some settings, MMSE-SIC shows stronger robustness than expected from synchronous comparisons. Overall, the thesis suggests that low-complexity SIC-based receivers can be made competitive when they are supported by multiple resources, soft information, variance-aware filtering, and receiver designs that reflect the actual uplink impairment model.

\section{Future Work}

Several research directions remain open:
\begin{itemize}
    \item Extend the simulation study to more users, more resources, and higher-order modulation such as QPSK or 16-QAM.
    \item Study fast-fading coded blocks, where a priori update iterations may provide larger gains than in block-fading channels.
    \item Develop lower-complexity approximations to MAP-IDD for HDS-CDMA and compare them with MMSE-SIC-IDD under the same latency constraints.
    \item Analyze imperfect CSI more deeply, including its effect on SIC ordering, cancellation accuracy, and variance-aware filtering.
    \item Optimize user pairing and resource assignment to balance SIC performance and user fairness.
    \item Extend the asynchronous OFDM study to random access scenarios with unknown timing offsets and blind or semi-blind timing estimation.
\end{itemize}
